% \section{TODO}
% \begin{itemize}

% \item \hl{supplement} Using
% MDS, we find that the subjects can be embedded as 20-dimensional vectors with less than 5
% in their pairwise distance (supplemental figure). 

% \item \hl{supplement} Indeed, a 2-dimensional projection given by PCA captures 72
% embedding (supplemental figure). 

% \item \hl{supplement} There
% are in fact a variety of supervised learning methods that enable us to make predictions on data in metric spaces
% (see supplemental text \\


% \item \hl{supplement} MDS of the
% dissimilarity matrix (sup fig) reveals that tuning width drives across-animals differences. \\

% \end{itemize}


\section{Introduction}


Measuring individual differences within and across species is fundamental to many areas of biology~\parencite{nelson1970outline,wagner1989biological}.
While such comparisons have yielded significant insights in neuroscience—particularly in surveys of gross-level neuroanatomy across species~\parencite{butler2005comparative} and in studies of small, well-characterized central pattern generator circuits ~\parencite{Marder2006}—it remains a major challenge to quantitatively compare the complex, high-dimensional dynamics of large neural systems, such as those found in mammalian brains.
REDUCE THIS TO BRING UP CONTROVERSIES
Over the past two decades, numerous methods have been developed to compare neural population dynamics. Pioneering work introduced Representational Similarity Analysis (RSA; \cite{Kriegeskorte2008_frontiers}), inspired by concepts from cognitive psychology \parencite{edelman1998representation}. Since then, several variants of RSA have been proposed \parencite[see][for a review]{kriegeskorte2019peeling}, alongside alternative approaches such as Canonical Correlations Analysis \parencite[CCA;][]{Raghu2017,gallego2020long} and its generalizations \parencite[reviewed in][]{zhuang2020technical}, Procrustes alignment \parencite{degenhart2020stabilization}, and hyperalignment \parencite{haxby2011common, haxby2020hyperalignment}. Similar techniques have been applied in studies of artificial deep networks, notably Centered Kernel Alignment \parencite[CKA;][]{Kornblith2019}, which was later found to be closely related to RSA~\parencite{Williams2024}. Additional approaches include linear classification to compare representational geometries across brain regions (e.g., \cite{bernardi}) and linear regression to quantify similarity between artificial and biological network representations (e.g., \cite{conwell2022can}). As a result, practitioners face a complex and interrelated ecosystem of quantitative methods for comparing representational geometries, and it remains unclear how these methods relate to one another or whether they can be unified within a common framework.

% Furthermore, few papers rigorously contend with the reality that these measures must be estimated from noisy neuronal responses over a limited number of trials \parencite[but see:][]{Cai2019,Schutt2021}.
% As a result, it is often unclear how to perform basic statistical practices such as a formal power analysis to estimate the number of trials and experimental conditions to include in a future study.


Crucially, scaling up these analyses to compare large collections of neural systems---across subjects, sessions, and brain regions (\textit{Fig.}\ref{fig:schematic}a-c)---remains a challenge.
For example, the International Brain Laboratory (IBL) released a dataset consisting of 121 recording sessions in 78 mice performing a visual discrimination task \parencite{iblreproducibility}.
The methodologies summarized above can quantify similarity between any pair of sessions (e.g. by computing the average canonical correlation coefficient).
Applied to the IBL dataset, this yields 7260 pairwise numbers.
How should researchers make sense of them?
One possibility is to show that they are higher, in aggregate, than those obtained under a null control (constructed, e.g., by shuffling the data).
But such aggregate tests do not chart a path towards more detailed questions. For instance, one might ask whether the recordings fall into clusters, e.g. small groups of subjects, or of brain regions, whose population activity sits closer to one another's than to the rest. An affirmative answer would establish that variability across recordings is structured rather than purely idiosyncratic. One might also ask whether the arrangement of these systems relative to one another supports prediction: whether a subject's behavioral profile can be inferred from other subjects with similar neural codes, or whether its neural activity can be decoded using a model fit entirely on different animals. Neither question can be answered from a list of pairwise numbers. Both presuppose that the systems occupy a common space, each recording holding a position within it and similarity acting as a distance.

STOP CONTROVERSIES HERE


% These questions are now pressing: recordings routinely span dozens of regions and large cohorts of subjects, yet whether population codes differ in any structured way across regions or individuals is far from settled—most task variables can be decoded almost anywhere in the brain, which has been read as evidence that representations are largely undifferentiated.

Here, we leverage a principled approach to geometrically compare activity measurements among multiple neural populations, such as many different brain regions or animal subjects, which is broadly referred to as \emph{shape distance}. 
Through analyses of simulated and experimental datasets we highlight two major advantages of this approach.
First, the shape distance between a pair of systems is quantified in terms of a transparently specified class of neural alignments (e.g. rotations + reflections, linear, or nonlinear transformations; \textit{Fig.}\ref{fig:schematic}e), thereby providing a flexible and interpretable family of measures.
Second, each distance within the shape metrics framework satisfies two key mathematical properties---symmetry and the triangle inequality---and therefore defines a metric space. This is precisely the setting the questions above require: an entire collection of recordings becomes a cloud of points (\textit{Fig.} \ref{fig:schematic}f), each neural system holding a position relative to every other, over which clustering and regression can be performed rigorously (\textit{Fig.}  \ref{fig:schematic}g). The mathematical foundations of the shape metrics framework were laid out in a series of recent conference papers \parencite{Williams2021,pospisil2023estimating,harvey2024duality,khosla2024soft}.
The contribution of the present manuscript is biological rather than mathematical: applying this framework across many brain regions and individuals in rodents and nonhuman primates, we show that neural population geometry is structured across both regions and individuals.

In the mouse postsubiculum, we uncover systematic animal-to-animal differences in head direction tuning statistics, trace these differences to tuning width via regression in shape space, and show that all subjects nonetheless share a common ring-shaped manifold, enabling decoding of head direction in a held-out animal from other animals' data alone. Applying the same framework across dozens of cortical, thalamic, and prefrontal regions in mice and non-human primates, we recover a hierarchical organization of brain areas in a fully unsupervised manner, notable given the growing appreciation that individual task variables are often decodable nearly everywhere in the brain \parencite{findling2025brain,angelaki2025brain, Lloreda2025,rosen2026distributed}. Finally, in a large cohort of mice performing a decision-making task, we show that individual differences in behavior are mirrored by individual differences in neural population geometry in some regions, allowing one animal's behavior to be predicted from the neural data of its closest neighbors. 

Together, these results show that neural population geometry is structured across both regions and individuals, organized by anatomy across brain areas and by behavior and genotype across individuals.